In late 2023 and early 2024, the primary way that the community evaluated code LMs was 
via benchmarks such as HumanEval \citep{chen2021evaluating} and MBPP \citep{austin2021program}, 
which test the ability to generate short code snippets from natural language 
specifications. While HumanEval and MBPP capture code generation abilities on simple and fundamental tasks, there was an absence of benchmarks capturing other 
fundamental dimensions of code LMs such as code understanding and execution. 

\begin{sloppypar}
Motivated by this, we introduced a new benchmark, \benchmark (\textbf{C}ode \textbf{R}easoning, \textbf{U}nderstanding, and e\textbf{X}ecution \textbf{Eval}uation) with two tasks: 1) output prediction, \benchmarko to measure code execution following and 2) input prediction, \benchmarki to measure code reasoning and understanding. An example of a sample in \benchmark is shown in Listings \ref{lst:benchmark-example1} and \ref{lst:benchmark-example} (modified for readability). \benchmark examines the abilities of code LMs to reason about the execution behavior of \textit{simple} Python programs. While LMs should not be expected to replace an interpreter on arbitrarily complex problems, we ensure the samples in our benchmark are simple (maximum 13 lines, no complex arithmetic) and solvable by a university-level CS graduate without needing extra memory. \benchmark provides a useful and important probe for better understanding the capabilities of code LMs, as following a few simple steps of code execution should be a basic requirement for these models. The ability to reason about the execution behavior of code also paves the way to tackling more difficult tasks such as code repair with execution feedback and code summarization. 
\end{sloppypar}
\noindent\begin{minipage}[t]{.48\textwidth}
\vspace{0pt}
\begin{lstlisting}[caption={Sample problem},label={lst:benchmark-example1}, captionpos=t, breaklines=true, language=Python, basicstyle=\ttfamily\scriptsize, breakindent=0pt, postbreak={}]
def f(string):
    string_x = string.rstrip("a")
    string = string_x.rstrip("e")
    return string

# output prediction, CRUXEval-O
assert f("xxxxaaee") == ??
## GPT4: "xxxx", incorrect

# input prediction, CRUXEval-I
assert f(??) == "xxxxaa"
## GPT4: "xxxxaae", correct
\end{lstlisting}
\end{minipage}\hfill
\begin{minipage}[t]{.48\textwidth}
\vspace{0pt}

\begin{lstlisting}[caption={Sample problem},label={lst:benchmark-example}, captionpos=t, breaklines=true, language=python, basicstyle=\ttfamily\scriptsize, breakindent=0pt, postbreak={}]
def f(nums):
    count = len(nums)
    for i in range(-count+1, 0):
        nums.append(nums[i])
    return nums
# output prediction, CRUXEval-O
assert f([2, 6, 1, 3, 1]) == ??
# GPT4: [2, 6, 1, 3, 1, 6, 1, 3, 1], incorrect

# input prediction, CRUXEval-I
assert f(??) == [2, 6, 1, 3, 1, 6, 3, 6, 6]
# GPT4: [2, 6, 1], incorrect
\end{lstlisting}
\end{minipage}

At a high level, our benchmark is constructed as follows. First, we use \codellamalarge to generate a large set of functions and inputs. The outputs are generated by executing the functions on the inputs. Second, we filter the set so that our benchmark only consists of short problems with low computation and memory requirements, problems which a good human programmer should be able to do without extra memory in a minute or so. Third, we randomly select 800 samples passing the filter, ensuring the benchmark is both small enough to easily run but large enough to reliably see performance differences among various models. We use this approach because while it is difficult to manually come up with examples where the strongest evaluated models fail completely, we observe that they fail quite often on random yet reasonable programs. We also highlight that as models improve, this generate-and-filter approach can be used to create future benchmarks that are more difficult and test different aspects of program execution.

Among the evaluated models, GPT-4 achieves a pass@1 of 67\% on \benchmarki and 63\% on \benchmarko. In contrast, the best evaluated open-source models only achieve 47\% on \benchmarki and 44\% on \benchmarko, failing over half the time at simple execution prediction and code reasoning despite being trained on 100G of Python code and 1T of code data. We also discover that despite its impressive abilities, GPT-4 consistently fails to understand the execution behavior of some surprisingly simple Python programs. 


%
% \if 0
% * We are interested in testing the extent to which success in existing benchmarks correlates with the ability of the models to execute simple programs. 

% * Our first contribution is a new benchmark suite designed to test execution ability. Focus on short programs that don't require extensive computation or arithmetic with large numbers. Our benchmark actually tests both the ability to execute forward and backward.

% * Using this benchmark, we show that there is a correlation between a models performance on HumanEval and its ability to model the execution of these short programs, but the correlation breaks down in some cases, particularly for small distilled models. 

% * Our experiments also show that both CoT and fine tuning can help improve the performance over the baseline, but the task remains hard even after using both of these techniques. 

% Overall, our benchmarks reveals that the gap which has been documented before between the highest performing models and the rest is not just limited to the existing benchmarks, but actually reflects a significantly stronger ability to reason about the behavior of code. Yet, despite their impressive abilities, even the best models still fail to predict the behavior of even surprisingly trivial programs. 

% \fi


% In Sec. \ref{sec:benchmark-construction}, we introduce our new benchmark of 800 Python functions, inputs, and outputs called \benchmark split into two subtasks, \benchmarki (input prediction) and \benchmarko (output prediction). In \benchmarko, the model is asked to execute a function on a given input, whereas in \benchmarki, is asked to find an input such that executing the function on the input gives a desired output. Alongside the benchmark, we include our three-step procedure used to generate the benchmark: large-scale distillation, filtering, and careful data size selection.

% In Sec. \ref{sec:evaluation}, \ref{sec:quantitative-analysis}, and Appendix \ref{sec:appendix-quantitative}, we evaluate a variety of models (Code Llama, StarCoder, WizardCoder, Phi, Phind, DeepSeek, Mistral, GPT). Pass@1 scores range from 14\% to 67\% on \benchmarki and 23\% to 63\% on \benchmarko. First, we discover that WizardCoder, Phind, and Phi have a HumanEval score disproportionately larger than their \benchmark score, suggesting that fine-tuning on GPT distilled data for code completion may not translate to general code reasoning and execution abilities. Second, we see a strong correlation between performance on the two tasks of our benchmark, even though they seem to be relatively different tasks.

% In Sec. \ref{subsec:cot}, we analyze the impact of using CoT on \benchmark. First, we see that GPT-4, the best performing model, enjoys a modest boost from CoT (67\% to 75\% for \benchmarki, 63\% to 81\% for \benchmarko). Second, we discover that CoT is generally more effective for output prediction than input prediction. Third, we see that CoT increase the diversity of generations. Finally, for Code Llama 13B, 34B, and GPT-3.5 (but not GPT-4), we find a surprisingly large number of samples for which the pass@1 with CoT is worse than that without.

% In Sec. \ref{subsec:finetuning}, we do a preliminary analysis of fine-tuning on input-output assertions on two different fine-tuning sets. In the first set, we decontaminate the fine-tuning set on problems in our benchmark and observe a modest improvement in accuracy: 41\% $\rightarrow$ 59\% on \benchmarko and 47\% $\rightarrow$ 63\% on \benchmarki. In the second set, we additionally include functions in the benchmark but with input-output pairs not in the benchmark. We see the accuracy further increase from 59\% $\rightarrow$ 63\% on \benchmarko and 63\% $\rightarrow$ 68\% on \benchmarki. Overall, while this simple fine-tuning scheme significantly helps benchmark performance, it is still far from acing the benchmark.

% Finally, in Sec. \ref{sec:qualitative-analysis} and Appendix \ref{sec:appendix-qualitative}), we present a qualitative analysis of successes and failures of GPT-4 on both tasks of our benchmark, revealing insights about GPT-4's behaviour on code execution reasoning. 

\subsection{Benchmark Construction} \label{sec:benchmark-construction}

\benchmark consists of 800 distinct functions, each with an input-output pair such that executing the function on the input deterministically produces the output. Using these functions and input-output pairs, we derive two benchmark tasks. In the \textit{output prediction} task, the goal is to predict the output of executing the function on its associated input. In the \textit{input prediction} task, the goal is to find any input such that executing the function on that input produces the output. For both tasks, we use an execution-based correctness metric. For input prediction, a generated input passes if \texttt{assert f(generated\_input) == output} passes, and for output prediction, a generated output passes if \texttt{assert f(input) == generated\_output} passes.

\subsubsection{Generating Candidates} \label{subsec:generating-candidates}
We use \codellamalarge to generate all the candidate functions and inputs of \benchmark. To do so, we provide the name of a function in the Python standard library such as \texttt{str.zfill} and ask for a Python function that uses the library function together with five test inputs. We use two varying few-shot examples during generation to improve diversity.

We use a total of 69 different functions from the standard library: 47 from the \texttt{str}, 11 from \texttt{dict}, and 11 from \texttt{list}. Overall, we generate a total of 102000 functions (46\% \texttt{str}, 27\% \texttt{dict}, 27\% \texttt{list}) and 489306 input-output pairs before filtering.

% String method count:  47 47000
% Dict method coun  11 27500
% List method count:  11 27500
% 102000

\subsubsection{Filtering Candidates} \label{subsec:filtering-candidates}
Next, we filter the generated candidates to ensure that the samples in the dataset are reasonable and of high quality. In order to avoid forcing the model to perform tasks such as arithmetic calculation, we design filters so that the benchmark only consists of samples that are solvable by a human without extra memory.

Concretely, we filter based on the following criteria. 
\begin{itemize}
    \item Compile time: all arguments of the function must be used in the function, length of code is between 75 and 300 characters, no syntax errors, proper assertion \texttt{assert f(input) == output}.

    \item Runtime: no float point operations, true division, exp, other integer operations must have at least one argument $\leq 3$, string and list operations must have at least one argument with length $\leq 3$, finish running in 2 seconds, no uncaught exceptions.

    \item Best effort to remove other undesirable code: function cannot have any imports (such as os, random), must be deterministic (random, set ordering), and cannot have side effects such as input, \texttt{\_\_builtins\_\_}.
\end{itemize}

\subsubsection{Data size and measuring noise} \label{subsec:data-size-noise}
The success of HumanEval (164 examples) shows that evaluation benchmarks can be small where faster and cheaper evaluation is an overlooked advantage. Since additional examples are easy to generate, we first overgenerate and then measure if the noise is sufficiently small on a smaller dataset. 

% When determining the final size of our benchmark, we balanced two factors: the benchmark should be small enough to easily run and iterate on but large enough so that signal overcomes noise. 

Out of all the samples, \codellamalarge outperforms \codellamamed as expected, and we would like to retain this property with high confidence in a smaller dataset. To do this, we took bootstrap samples of size $N$ out of $\sim$1700 samples to measure the probability that the performance would be reversed, shown in Fig. \ref{fig:benchmark-filtering-size-selection}. We found that 800 examples are enough to test that \codellamalarge $>$ \codellamamed, \textsc{Code Llama cot} $>$ \textsc{Code Llama}, and \textsc{DeepSeek 33B} $>$ \codellamalarge on output prediction.


We measure two sources of noise: 1) sampling which data points to include in the benchmark, and 2) sampling candidates from models for each data point (temperature $>0$). Of these, 1) dominates 2). For 1) since model $A$ does not always outperform model $B$ on all data points even if $A>B$ in aggregate, the measured performance depends on which data points are included. We can measure both noise on each model individually, and also measure type 1) noise on pairs of models using bootstrap. Fortunately, we do not see major differences between models and the most important factor is just the size of dataset. Type 1) noise is generally around 1.5\% for each model whereas type 2) is around 0.2\% at $N=800$. Type 1) noise usually becomes smaller on pairs of models due to correlation, yielding statistically significant results at the $\alpha=0.05$ level for many model pairs.

\begin{figure}[H]
    \centering
    \includegraphics[width=0.49\textwidth] {figs/wrong_order/wrong_order_box_input_pass1.pdf}
    \includegraphics[width=0.49\textwidth]{figs/wrong_order/wrong_order_box_output_pass1.pdf}
    \caption{Difference between model pairs on bootstrap samples of various sizes. The whiskers show (2.5, 97.5) and boxes show (25, 75) percentiles. }
    \label{fig:benchmark-filtering-size-selection}
\end{figure}

\subsection{Evaluation} \label{sec:evaluation}

We evaluate a selection of models on \benchmark: StarCoder (Base 7B, 15.5B) \citep{li2023starcoder}, Mistral (7B) \citep{jiang2023mistral}, WizardCoder (13B, 34B) \citep{luo2023wizardcoder}, Phi-1 \cite{gunasekar2023textbooks} and Phi-1.5 \citep{li2023textbooks} (1.3B), Phind v2 \citep{Phind} (34B), Code Llama \citep{roziere2023code} (Base and Python 7B, 13B, 34B), DeepSeek Coder (Base and Instruct 6.7B, 33B), GPT-3.5 \citep{brown2020language, ouyang2022training}, and GPT-4 \citep{openai2023gpt}.

We use $N=100$ samples for all non-GPT models and $N=10$ samples for GPT models. We report both pass@1 scores $(T=0.2)$ and pass@5 scores $(T=0.8)$. The results are shown in Fig.~\ref{fig:main-results}. In Fig. \ref{fig:main-results}, we show the intervals generated by 10000 bootstrap samples from the dataset, where non-overlapping whiskers would be significant at the 2.5\% level. To get more statistical power, we also compare pairs of models on each bootstrapped sample.
The intervals generally decreases due to correlations. On all models vs. \codellamalarge, if the median bar clears the whisker in Fig.~\ref{fig:main-results}, then the difference actually holds with $>$97.5\% probability under paired bootstrap. For example, \codellamalarge is better than \textsc{wizard\_34B} on input and \codellamalarge is worse than \textsc{deepseek\_33B} on output prediction with $>$97.5\% probability.

\begin{figure}[H]
     \centering
     \begin{subfigure}[b]{0.49\textwidth}
         \centering
         \includegraphics[width=\textwidth]{figs/main_results/main_box_input.pdf}
         \subcaption{\benchmarki Performance}
         \label{fig:main-results-pass1-input}
     \end{subfigure}
     \hfill
     \begin{subfigure}[b]{0.49\textwidth}
         \centering
         \includegraphics[width=\textwidth]{figs/main_results/main_box_output.pdf}
         \subcaption{\benchmarko Performance}
         \label{fig:main-results-pass1-output}
     \end{subfigure}
     \caption{Main Results. Boxes show (25, 75) percentiles, whiskers show (2.5, 97.5), and the middle bar shows the median ($\approx$ mean). }
     \label{fig:main-results}
\end{figure}


\subsection{Impact of Anonymizing Functions} \label{sec:cruxeval-anonymization}
One potential source of bias is that models may have memorized the functions in our benchmark.
While we believe this is not the case due to the synthetically generated nature of the benchmark, we 
conduct another experiment where we replace the variable names of functions with synthetic "anonymous" 
identifiers while keeping function semantics intact. We test CodeLlama 7B, 13B, and 34B on a subset of our benchmark
where variable names are replaced with \texttt{x1, x2, ...} identifiers. An example of an anonymized function is shown 
in Listing \ref{lst:benchmark-anonymized-sample}. We use the same few-shot prompt 
without anonymization and report both pass@1 ($T=0.2$) and pass@5 ($T=0.8$) results 
on the anonymized benchmark with $N=10$ samples. The results are shown in 
Table \ref{tab:benchmark-results-anonymous}. This strengthens the case against 
memorization effects.

\begin{lstlisting}[caption={Sample of benchmark and anonymized version},label={lst:benchmark-anonymized-sample}, captionpos=t, breaklines=true]
Original:
def f(s):
    nums = ''.join(filter(lambda c:c.isdecimal(), s))
    if nums == '': return 'none'
    m = max([int(num) for num in nums.split(',')])
    return str(m)
assert f('01,001') == '1001'

Anonymized:
def f(x0):
    x1 = ''.join(filter(lambda x2: x2.isdecimal(), x0))
    if x1 == '':
        return 'none'
    x3 = max([int(x4) for x4 in x1.split(',')])
    return str(x3)
assert f('01,001') == '1001'
\end{lstlisting}

\begin{table}[htbp]
    \centering
    \caption{Impact of Anonymization on \benchmark}
    \begin{tabular}{cccccc}
        \toprule
        \multirow{2}{*}{\textbf{Model}} & \multirow{2}{*}{\textbf{Anonymized}} & \multicolumn{2}{c}{\textbf{Input Prediction}} & \multicolumn{2}{c}{\textbf{Output Prediction}} \\
        \cmidrule(lr){3-4} \cmidrule(lr){5-6}
        & & \textbf{Pass@1} & \textbf{Pass@5} & \textbf{Pass@1} & \textbf{Pass@5} \\
        \midrule
        \multirow{3}{*}{CodeLlama 7B} 
	& \xmark & 36.6\% & 48.0\% & 36.4\% & 43.5\% \\
        & \cmark & 37.5\% & 53.3\% & 34.0\% & 46.9\% \\
        & $\Delta$ & \green{+0.9\%} & \green{+5.3\%} & \red{-2.4\%} & \green{+3.4\%} \\
        \midrule
        \multirow{3}{*}{CodeLlama 13B} 
        & \xmark & 39.0\% & 50.2\% & 38.3\% & 44.7\% \\
        & \cmark & 40.0\% & 55.8\% & 36.1\% & 50.6\% \\
        & $\Delta$ & \green{+1.0\%} & \green{+5.6\%} & \red{-2.2\%} & \green{+5.9\%} \\
        \midrule
        \multirow{3}{*}{CodeLlama 34B} 
        & \xmark & 46.5\% & 57.4\% & 41.1\% & 47.5\% \\
        & \cmark & 48.0\% & 63.8\% & 39.1\% & 54.0\% \\
        & $\Delta$ & \green{+1.5\%} & \green{+6.4\%} & \red{-2.0\%} & \green{+6.5\%} \\
        \bottomrule
    \end{tabular}
    \label{tab:benchmark-results-anonymous}
\end{table}

\subsection{Discussion}
In CRUXEval, we develop a new benchmark and metric to evaluate code reasoning and execution through two
related tasks: input prediction and output prediction. Our experiments provide evidence that 
models do indeed have some ability to perform these tasks. However, even the best models still
showed significant room for improvement, considering that these tasks only required basic
program reasoning abilities. As we discuss later in this thesis, CRUXEval was used as a 
core benchmark by the community to track progress on code reasoning and execution, inspiring
approaches such as training on execution traces in addition to code.
